\documentclass[../../main.tex]{subfiles}
\begin{document}
\section{Riemann 体积形式}\label{section:ch15-volume-form}

设 $(M,g)$ 是一个正维数的有向 Riemann 流形。由命题\ref{proposition:ch13-gram-schmidt}，我们知道在 $M$ 每一点的某个邻域上存在一个光滑标准正交标架 $(E_1,\dots,E_n)$。必要时把 $E_1$ 换成 $-E_1$，便可在每一点的某个邻域上找到一个定向标准正交标架。

\begin{proposition}\label{proposition:ch15-volume-form}
设 $(M,g)$ 是一个带边或不带边的有向 Riemann $n$ 流形，$n\geqslant1$。存在唯一一个光滑定向形式 $\omega_g\in\Omega^n(M)$，称为\emph{Riemann 体积形式}，它对 $M$ 的每个局部定向标准正交标架 $(E_i)$ 都满足
\begin{align}
\omega_g(E_1,\dots,E_n)=1.\label{align:ch15-volume-normal}
\end{align}
\end{proposition}

\begin{proof}
首先假设这样的形式 $\omega_g$ 存在。若 $(E_1,\dots,E_n)$ 是开子集 $U\subseteq M$ 上的任意局部定向标准正交标架，$(\varepsilon^1,\dots,\varepsilon^n)$ 是对偶余标架，则在 $U$ 上可写 $\omega_g=f\,\varepsilon^1\wedge\cdots\wedge\varepsilon^n$。条件\eqref{align:ch15-volume-normal} 此时化为 $f=1$，故
\begin{align}
\omega_g=\varepsilon^1\wedge\cdots\wedge\varepsilon^n.\label{align:ch15-volume-orthonormal}
\end{align}
这就证明了这样的形式唯一确定。

为证存在性，我们希望在每一点的邻域上由\eqref{align:ch15-volume-orthonormal} 定义 $\omega_g$，故需检查这个定义与定向标准正交标架的选取无关。若 $(\widetilde E_1,\dots,\widetilde E_n)$ 是另一个定向标准正交标架，对偶余标架为 $(\widetilde\varepsilon^1,\dots,\widetilde\varepsilon^n)$，令 $\widetilde\omega_g=\widetilde\varepsilon^1\wedge\cdots\wedge\widetilde\varepsilon^n$。我们可写 $\widetilde E_i=A^j_iE_j$，其中 $(A^j_i)$ 是某个光滑函数矩阵。两个标架都标准正交这一事实意味着 $(A^j_i(p))\in O(n)$，故 $\det(A^j_i)=\pm1$；而两个标架一致定向这一事实迫使符号为正。我们计算
\begin{align*}
\omega_g(\widetilde E_1,\dots,\widetilde E_n)
=\det\bigl(\varepsilon^j(\widetilde E_i)\bigr)
=\det(A^j_i)=1=\widetilde\omega_g(\widetilde E_1,\dots,\widetilde E_n).
\end{align*}
因此 $\omega_g=\widetilde\omega_g$，故在每一点的邻上关于某个光滑定向标准正交标架由\eqref{align:ch15-volume-orthonormal} 定义 $\omega_g$，便得到一个整体 $n$-形式。所得形式显然光滑，且对每个定向标准正交标架满足\eqref{align:ch15-volume-normal}。
\end{proof}

\begin{exercise}\label{exercise:ch15-30}
设 $(M,g)$ 与 $(\widetilde M,\widetilde g)$ 是带边或不带边的正维 Riemann 流形，$F:M\to\widetilde M$ 是一个局部等距。证明 $F^*\omega_{\widetilde g}=\omega_g$。
\end{exercise}

虽然 Riemann 体积形式关于一个定向标准正交标架的表达式特别简单，但在坐标下给出它的表达式也是有用的。

\begin{proposition}\label{proposition:ch15-volume-coords}
设 $(M,g)$ 是一个带边或不带边的有向 Riemann $n$ 流形，$n\geqslant1$。在任意定向光滑坐标 $(x^i)$ 下，Riemann 体积形式有局部坐标表达式
\begin{align*}
\omega_g=\sqrt{\det(g_{ij})}\,\mathrm{d}x^1\wedge\cdots\wedge\mathrm{d}x^n,
\end{align*}
其中 $g_{ij}$ 是 $g$ 在这些坐标下的分量。
\end{proposition}

\begin{proof}
设 $(U,(x^i))$ 是一个定向光滑坐标图，$p\in M$。在这些坐标下，$\omega_g=f\,\mathrm{d}x^1\wedge\cdots\wedge\mathrm{d}x^n$，其中 $f$ 是某个正的系数函数。为计算 $f$，取在 $p$ 的某个邻域上定义的任意光滑定向标准正交标架 $(E_i)$，令 $(\varepsilon^i)$ 为对偶余标架。若我们把坐标标架用标准正交标架写成
\begin{align*}
\frac{\partial}{\partial x^i}=A^j_iE_j,
\end{align*}
则可计算
\begin{align*}
f=\omega_g\left(\frac{\partial}{\partial x^1},\dots,\frac{\partial}{\partial x^n}\right)
=\varepsilon^1\wedge\cdots\wedge\varepsilon^n\left(\frac{\partial}{\partial x^1},\dots,\frac{\partial}{\partial x^n}\right)
=\det\left(\varepsilon^j\left(\frac{\partial}{\partial x^i}\right)\right)
=\det(A^j_i).
\end{align*}
另一方面，注意
\begin{align*}
g_{ij}=\left\langle\frac{\partial}{\partial x^i},\frac{\partial}{\partial x^j}\right\rangle_g
=\left\langle A^k_iE_k,A^l_jE_l\right\rangle_g
=A^k_iA^l_j\langle E_k,E_l\rangle_g
=\sum_k A^k_iA^k_j.
\end{align*}
最后这个表达式是矩阵乘积 $A^{\mathsf T}A$ 的 $(i,j)$ 元，其中 $A=(A^j_i)$。因此
\begin{align*}
\det(g_{ij})=\det(A^{\mathsf T}A)=\det A^{\mathsf T}\det A=(\det A)^2,
\end{align*}
由此得 $f=\det A=\pm\sqrt{\det(g_{ij})}$。由于 $(\partial/\partial x^i)$ 与 $(E_j)$ 都定向，符号必为正。
\end{proof}

\subsection{Riemann 流形中的超曲面}

设 $(M,g)$ 是一个带边或不带边的有向 Riemann 流形，$S\subseteq M$ 是一个带边或不带边的浸入超曲面。沿 $S$ 的任意单位法向量场都处处不与 $S$ 相切，故由命题\ref{proposition:ch15-hypersurface-orientation} 它确定 $S$ 的一个定向。下一个命题给出关于这个定向、$S$ 上诱导度量的体积形式的一个简单公式。

\begin{proposition}\label{proposition:ch15-hypersurface-volume}
设 $(M,g)$ 是一个带边或不带边的有向 Riemann 流形，$S\subseteq M$ 是一个带边或不带边的浸入超曲面，$\widetilde g$ 是 $S$ 上的诱导度量。设 $N$ 是沿 $S$ 的一个光滑单位法向量场。关于由 $N$ 确定的 $S$ 的定向，$(\widetilde S,\widetilde g)$ 的体积形式由
\begin{align*}
\omega_{\widetilde g}=\iota_S^*(N\lrcorner\omega_g)
\end{align*}
给出。
\end{proposition}

\begin{proof}
由命题\ref{proposition:ch15-hypersurface-orientation}，$(n-1)$-形式 $\iota_S^*(N\lrcorner\omega_g)$ 是 $S$ 的一个定向形式。为证它是诱导 Riemann 度量的体积形式，我们只需证明：只要它作用于 $S$ 的一个定向标准正交标架，就取值 $1$。于是设 $(E_1,\dots,E_{n-1})$ 是这样一个标架。在每个 $p\in S$ 处，基 $(N_p,E_1|_p,\dots,E_{n-1}|_p)$ 标准正交（图\ref{figure:ch15-hypersurface}），并且对 $T_pM$ 是定向的（这正是由 $N$ 确定的定向的定义）。因此
\begin{align*}
\iota_S^*(N\lrcorner\omega_g)(E_1,\dots,E_{n-1})
=\omega_g(N,E_1,\dots,E_{n-1})=1,
\end{align*}
这就证明了结果。
\end{proof}

\begin{figure}[htbp]
\centering
\includegraphics[width=0.8\textwidth]{fig15-8-hypersurface.png}
\caption{Riemann 流形中的一个超曲面}
\label{figure:ch15-hypersurface}
\end{figure}

由于下面的命题，命题\ref{proposition:ch15-hypersurface-volume} 的结果在带边 Riemann 流形的情形中尤为重要。

\begin{proposition}\label{proposition:ch15-boundary-normal}
设 $M$ 是任意一个带边 Riemann 流形。沿 $\partial M$ 存在唯一一个光滑向外单位法向量场 $N$。
\end{proposition}

\begin{proof}
首先证明唯一性。在任意点 $p\in\partial M$ 处，子空间 $(T_p\partial M)^\perp\subseteq T_pM$ 是一维的，故在 $p$ 处恰有两个与 $\partial M$ 正交的单位向量。由于任意单位法向量 $N$ 处处不与 $\partial M$ 相切，它在任意光滑边界坐标图中必有非零的 $x^n$ 分量。因此，两个单位法向量选择中恰有一个的 $x^n$ 分量为负，这等价于向外指向。

为证存在性，令 $f:M\to\mathbb{R}$ 是一个边界定义函数（命题\ref{proposition:ch05-boundary-defining}），令 $N$ 为单位向量场 $-\operatorname{grad}f/|\operatorname{grad}f|_g$ 在 $\partial M$ 上的限制。由于在 $\partial M$ 各点处 $\mathrm{d}f\ne0$，$N$ 在 $\partial M$ 上良定义且光滑。于是由习题\ref{exercise:ch13-21}，$N$ 与 $\partial M$ 正交；由习题\ref{exercise:ch05-44}，它向外指向，因为
\begin{align*}
Nf=-\frac{\langle\operatorname{grad}f,\operatorname{grad}f\rangle_g}{|\operatorname{grad}f|_g}=-|\operatorname{grad}f|_g<0.
\end{align*}
\end{proof}

下一个推论是直接的。

\begin{corollary}\label{corollary:ch15-boundary-volume}
若 $(M,g)$ 是一个有向带边 Riemann 流形，$\widetilde g$ 是 $\partial M$ 上的诱导 Riemann 度量，则 $\widetilde g$ 的体积形式为
\begin{align*}
\omega_{\widetilde g}=\iota_{\partial M}^*(N\lrcorner\omega_g),
\end{align*}
其中 $N$ 是沿 $\partial M$ 的向外单位法向量场。
\end{corollary}
\end{document}
